\documentclass[11pt]{article}

\usepackage[margin=1in]{geometry}
\usepackage{amsmath, amssymb, amsthm}
\usepackage{mathtools}
\usepackage{bm}
\usepackage{xcolor}
\usepackage{hyperref}
\hypersetup{colorlinks=true, linkcolor=blue!50!black, citecolor=blue!50!black}

\newtheorem{theorem}{Theorem}
\newtheorem{lemma}[theorem]{Lemma}
\newtheorem{proposition}[theorem]{Proposition}
\newtheorem{definition}[theorem]{Definition}
\newtheorem{remark}[theorem]{Remark}

\DeclareMathOperator{\Tr}{Tr}
\DeclareMathOperator{\tr}{tr}
\DeclareMathOperator{\vecop}{vec}
\newcommand{\IR}{\mathbb{R}}
\newcommand{\IE}{\mathbb{E}}
\newcommand{\bP}{\boldsymbol{P}}
\newcommand{\bbeta}{\boldsymbol{\beta}}
\newcommand{\Wtilde}{\widetilde{W}}
\newcommand{\utilde}{\widetilde{u}}

\title{A closed-form expression for $\Tr\!\bigl(M\,(A\otimes A)\bigr)$\\
when $A$ is given as a Matrix Product Operator}
\author{}
\date{}

\begin{document}
\maketitle

\section{Goal and statement}

Fix integers $m \ge 1$, $N \ge 3$, $d \ge 2$, set $D = d^N$, and write
$\bP = \{A, U_1, U_2, 4K\}$ for the four ``replica patterns'' of
$\{a,b,c,d\}$ defined in equation~\eqref{eq:patterns} below. Let
\[
M \in \IR^{D^2 \times D^2}
\]
be the second-moment matrix of the $TT$ sketch, with rows indexed by
$(i^a, i^c) \in [d]^N \times [d]^N$ and columns by $(i^b, i^d)$,
\begin{equation}\label{eq:Mdef}
M\bigl[(i^a,i^c),\,(i^b,i^d)\bigr] \;=\;
\IE\!\left[\,(S^\top S)_{i^a,i^b}\,(S^\top S)_{i^c,i^d}\,\right].
\end{equation}
Thanks to the chain construction of the thesis, $M$ admits a description
as a finite sum over pattern sequences $\bP^N$, with per-site pattern
transfer matrices $W^P$ and boundary vectors $u_{\text{init}}^{(P_1,P_2)}$,
$u_{\text{end}}^{P_N}$, all of which are explicit polynomials in
$(m, m_1, m_2, \sigma_2, \sigma_4)$ where $m=m_1m_2$ is the gadget
factorization and $(\sigma_2,\sigma_4)$ are the second and fourth moments
per entry. We recall this construction in
Section~\ref{sec:M-chain}.

Suppose now that we are given a symmetric matrix
$A \in \IR^{D \times D}$ in \emph{Matrix Product Operator (MPO) form},
that is, parameterized by $N$ tensors
\[
A^{(l)} \in \IR^{r_L^{(l)} \times d \times d \times r_R^{(l)}}, \qquad l=1,\dots,N,
\]
with $r_L^{(1)} = r_R^{(N)} = 1$ (rank-$1$ boundaries) and matching internal
bonds $r_R^{(l)} = r_L^{(l+1)} =: r_l$. The MPO defines the entries of $A$ via
\begin{equation}\label{eq:Adef}
A_{ij} \;=\; A^{(1)}_{i_1 j_1}\,A^{(2)}_{i_2 j_2} \cdots A^{(N)}_{i_N j_N},
\qquad i,j\in[d]^N,
\end{equation}
where $i = (i_1,\dots,i_N)$, $j = (j_1,\dots,j_N)$, and we write
$A^{(l)}_{ab} \in \IR^{r_{l-1} \times r_l}$ for the matrix obtained by
fixing the two physical indices to $(a,b)$. The right-hand side of
\eqref{eq:Adef} is therefore a product of $N$ matrices of sizes
$1\times r_1, r_1\times r_2,\dots, r_{N-1}\times 1$, which evaluates to a
scalar.

\medskip

The goal of this note is to prove and verify the following closed-form
expression for the bilinear form $\Tr\bigl(M(A\otimes A)\bigr)$.

\begin{theorem}[Closed form]\label{thm:main}
Let $M$ be the matrix defined in \eqref{eq:Mdef} and let $A$ be the matrix
defined in \eqref{eq:Adef}. For every site $l\in\{1,\dots,N\}$ and every
pattern $P\in\bP$, define the per-site $r_{l-1}^2 \times r_l^2$ matrix
\begin{equation}\label{eq:beta-def}
\beta_P^{(l)}\bigl[(\alpha,\gamma),(\alpha',\gamma')\bigr]
\;=\;
\sum_{(a,b,c,d)\in P}
A^{(l)}_{ab}[\alpha,\alpha']\;A^{(l)}_{cd}[\gamma,\gamma'],
\end{equation}
where the sum is over $4$-tuples $(a,b,c,d)\in[d]^4$ that realize the
pattern $P$ (Definition~\ref{def:patterns} below). Then
\begin{equation}\label{eq:thm-main}
\boxed{\;
\Tr\bigl(M\,(A\otimes A)\bigr) \;=\;
\utilde_{\mathrm{end}}^{\!\top}
\;\Wtilde^{(N-1)}\,\Wtilde^{(N-2)}\cdots\Wtilde^{(3)}\;
\utilde_{\mathrm{init}}\;}
\end{equation}
where the aggregated quantities are
\begin{align}
\Wtilde^{(l)}
&= \sum_{P\in\bP}\, W^P \,\otimes\, \bigl(\beta_P^{(l)}\bigr)^{\!\top}
&& (l = 3,\dots,N-1)\;,\label{eq:Wtilde}\\[2pt]
\utilde_{\mathrm{init}}
&= \sum_{P_1,P_2\in\bP}\,
u_{\mathrm{init}}^{(P_1,P_2)}\,\otimes\,
\vecop\!\bigl(\beta_{P_1}^{(1)}\,\beta_{P_2}^{(2)}\bigr),
\label{eq:uinit}\\[2pt]
\utilde_{\mathrm{end}}
&= \sum_{P\in\bP}\,
u_{\mathrm{end}}^{P}\,\otimes\,
\vecop\!\bigl(\beta_P^{(N)}\bigr).
\label{eq:uend}
\end{align}
\end{theorem}

The right-hand side of \eqref{eq:thm-main} is a length-$(N-2)$ matrix-vector
product on a space of dimension $4 r^2$ (with $r := \max_l r_l$), and is
therefore computable in $O(N r^6)$ floating-point operations. Both
$\utilde_{\mathrm{init}}$ and $\utilde_{\mathrm{end}}$ are vectors of length
$4 r^2$.

The remainder of the note proves \eqref{eq:thm-main}. We first introduce
notation (Section~\ref{sec:notation}), recall the chain structure of $M$
(Section~\ref{sec:M-chain}), establish the MPO-times-MPO factorization
of the joint contraction (Section~\ref{sec:mpo-double}), and finally
combine the two chains into the aggregated form
(Section~\ref{sec:assemble}). A numerical verification across six
configurations agrees with the brute-force evaluation
$\vecop(A)^\top M \,\vecop(A)$ to machine precision (Table~\ref{tab:verify}).

\section{Notation and conventions}\label{sec:notation}

\subsection{Indices and tensor product}
We use multi-indices $i = (i_1,\dots,i_N)\in[d]^N$ for the rows and columns
of operators on $\IR^D$. The Kronecker product of matrices
$X\in\IR^{m\times n}$ and $Y\in\IR^{p\times q}$ is the
$mp\times nq$ matrix $X\otimes Y$ with entries
\[
(X\otimes Y)\bigl[(i,k),\,(j,l)\bigr] \;=\; X_{ij}\,Y_{kl},
\]
where the row index $(i,k)$ is read as $i\cdot p + k$ and the column index
$(j,l)$ as $j\cdot q + l$. The vectorization $\vecop:\IR^{p\times q}\to\IR^{pq}$
is defined by $\vecop(X)_{ip+k} = X_{ik}$ (row-major). Two identities used
repeatedly are
\begin{equation}\label{eq:kron-identities}
(X\otimes Y)(Z\otimes W) \;=\; (XZ)\otimes(YW),
\qquad
\Tr(X\otimes Y) \;=\; \Tr(X)\,\Tr(Y).
\end{equation}

\subsection{The four replica patterns}\label{def:patterns}
A \emph{pattern} $P\in\bP$ is an equivalence relation on the set
$\{a,b,c,d\}$ of four ``replica labels'' such that every block has even
size. There are exactly four such relations: the trivial one and three
pairings, namely
\begin{equation}\label{eq:patterns}
\begin{aligned}
4K  &: \{a=b=c=d\}, \\
A   &: \{a=b\}\sqcup\{c=d\}, \\
U_1 &: \{a=c\}\sqcup\{b=d\}, \\
U_2 &: \{a=d\}\sqcup\{b=c\}.
\end{aligned}
\end{equation}
Given a $4$-tuple $(a,b,c,d)\in[d]^4$, we write $(a,b,c,d)\in P$ to mean
that the equality pattern of $(a,b,c,d)$ equals $P$. Explicitly,
\[
\begin{aligned}
(a,b,c,d)\in 4K  &\iff a=b=c=d,\\
(a,b,c,d)\in A   &\iff a=b,\;c=d,\;a\ne c,\\
(a,b,c,d)\in U_1 &\iff a=c,\;b=d,\;a\ne b,\\
(a,b,c,d)\in U_2 &\iff a=d,\;b=c,\;a\ne b.
\end{aligned}
\]
Tuples outside these four patterns have an odd-sized block and contribute
zero to all moments below; we ignore them.

\section{Chain structure of $M$}\label{sec:M-chain}

We recall the form of $M$ produced by the thesis construction. (No new
material; this section fixes notation that is used throughout.)

\begin{proposition}[Thesis chain decomposition]\label{prop:M-chain}
There exist scalar quantities $V(P_1,\dots,P_N)\in\IR$ indexed by pattern
sequences $\bP^N$ such that, for every choice of multi-indices
$i^a,i^b,i^c,i^d\in[d]^N$,
\begin{equation}\label{eq:M-as-sum}
M\bigl[(i^a,i^c),(i^b,i^d)\bigr]
\;=\;
\sum_{(P_1,\dots,P_N)\in\bP^N}
\Bigl(\,\textstyle\prod_{l=1}^N \mathbf{1}\!\bigl[(i^a_l,i^b_l,i^c_l,i^d_l)\in P_l\bigr]\Bigr)
\,V(P_1,\dots,P_N).
\end{equation}
Moreover, $V$ has the chain (transfer-matrix) factorization
\begin{equation}\label{eq:V-chain}
V(P_1,\dots,P_N)
\;=\;
\bigl(u_{\mathrm{end}}^{P_N}\bigr)^{\!\top}\,
W^{P_{N-1}}\, W^{P_{N-2}}\cdots W^{P_3}\,
u_{\mathrm{init}}^{(P_1,P_2)},
\end{equation}
where $W^P\in\IR^{4\times 4}$ for $P\in\bP$ and
$u_{\mathrm{init}}^{(P_1,P_2)}, u_{\mathrm{end}}^{P_N}\in\IR^4$
are the explicit polynomials in $(m,m_1,m_2,\sigma_2,\sigma_4)$ given
by formulas (19), (20), (21) of the thesis.
\end{proposition}

The interpretation of \eqref{eq:M-as-sum} is straightforward: $M$ is built
by summing, over all $4^N$ assignments of a per-site pattern, the product
of an indicator (only those $(i^a,i^b,i^c,i^d)$ that match the pattern
contribute) and a chain value $V$.

\section{Goal as an explicit four-fold sum}\label{sec:goal}

The bilinear form $\Tr(M\cdot Y)$ for any $Y\in\IR^{D^2\times D^2}$ is
\[
\Tr(M Y) \;=\; \sum_{I,J}\, M_{IJ}\,Y_{JI},
\]
where $I,J\in[D^2]$. With $Y = A\otimes A$, decoded into the user's
$(i^a,i^b,i^c,i^d)$ convention,
\[
(A\otimes A)\bigl[(i^b,i^d),(i^a,i^c)\bigr] \;=\;
A_{i^b i^a}\,A_{i^d i^c} \;=\; A_{i^a i^b}\,A_{i^c i^d}
\quad\text{(since $A=A^\top$)}.
\]
Hence
\begin{equation}\label{eq:trMB-base}
\Tr\bigl(M(A\otimes A)\bigr)
\;=\;
\sum_{i^a,i^b,i^c,i^d \in [d]^N}
A_{i^a i^b}\,A_{i^c i^d}\,
M\bigl[(i^a,i^c),(i^b,i^d)\bigr].
\end{equation}
This identity is independent of any sketch structure and holds for any
symmetric $A$ and any $M$ in the user's convention.

\section{The MPO factorization of $A_{i^a i^b}\,A_{i^c i^d}$}\label{sec:mpo-double}

Substituting~\eqref{eq:Adef} into the product yields, with double indices
on the bond legs $(\alpha_l)_{l=0}^N$ and $(\gamma_l)_{l=0}^N$
(both with closed boundaries $\alpha_0=\alpha_N=\gamma_0=\gamma_N=1$),
\[
A_{i^a i^b}\,A_{i^c i^d}
\;=\;
\Bigl(\prod_{l=1}^N
A^{(l)}_{i^a_l i^b_l}[\alpha_{l-1},\alpha_l]\Bigr)\,
\Bigl(\prod_{l=1}^N
A^{(l)}_{i^c_l i^d_l}[\gamma_{l-1},\gamma_l]\Bigr).
\]
Combining the two products site-by-site as a tensor product on the bond legs,
\begin{equation}\label{eq:AA-mpo-prod}
A_{i^a i^b}\,A_{i^c i^d}
\;=\;
\prod_{l=1}^N
\bigl(A^{(l)}_{i^a_l i^b_l} \otimes A^{(l)}_{i^c_l i^d_l}\bigr)\!
\bigl[(\alpha_{l-1},\gamma_{l-1}),\,(\alpha_l,\gamma_l)\bigr],
\end{equation}
where the product is now a matrix product on the joint bond space
$\IR^{r_{l-1}^2}\to\IR^{r_l^2}$, evaluated at the closed boundary indices.

\section{Combining the two chains}\label{sec:assemble}

We now substitute \eqref{eq:M-as-sum} and \eqref{eq:AA-mpo-prod}
into~\eqref{eq:trMB-base} and exchange sums.

\subsection{Per-site collapse via $\beta_P^{(l)}$}

Plugging \eqref{eq:M-as-sum} into \eqref{eq:trMB-base}:
\[
\Tr(M(A\otimes A))
\;=\;
\sum_{(P_1,\dots,P_N)}
V(P_1,\dots,P_N)
\sum_{\substack{i^a,i^b,i^c,i^d:\\ (i^a_l,i^b_l,i^c_l,i^d_l)\in P_l\,\forall l}}
A_{i^a i^b}\,A_{i^c i^d}.
\]
The inner sum factorizes site by site because the indicator constraints and
the MPO product~\eqref{eq:AA-mpo-prod} are both products over $l$:
\begin{align*}
&\sum_{\substack{i^a,i^b,i^c,i^d:\\ (i^a_l,i^b_l,i^c_l,i^d_l)\in P_l\,\forall l}}
A_{i^a i^b}\,A_{i^c i^d} \\
&\quad=\;
\sum_{\alpha,\gamma}
\prod_{l=1}^N\;
\sum_{(a,b,c,d)\in P_l}
\bigl(A^{(l)}_{ab}\otimes A^{(l)}_{cd}\bigr)\!
\bigl[(\alpha_{l-1},\gamma_{l-1}),(\alpha_l,\gamma_l)\bigr] \\
&\quad=\;
\sum_{\alpha,\gamma}
\prod_{l=1}^N \beta_{P_l}^{(l)}\bigl[(\alpha_{l-1},\gamma_{l-1}),(\alpha_l,\gamma_l)\bigr],
\end{align*}
where in the last step we used the definition~\eqref{eq:beta-def} of
$\beta_P^{(l)}$. The remaining sum over closed bond indices
$(\alpha,\gamma)$ is, by inspection, the $(1,1)$-entry of the matrix product
of $\beta$'s:
\begin{equation}\label{eq:beta-trace}
\sum_{\substack{i^a,i^b,i^c,i^d:\\ (i^a_l,i^b_l,i^c_l,i^d_l)\in P_l\,\forall l}}
A_{i^a i^b}\,A_{i^c i^d}
\;=\;
\bigl(\beta_{P_1}^{(1)}\,\beta_{P_2}^{(2)}\cdots\beta_{P_N}^{(N)}\bigr)\!\bigl[1,1\bigr].
\end{equation}
(Recall that $r_0 = r_N = 1$, so this product has shape $1\times 1$.)

\subsection{Joining the two chains}

Substituting \eqref{eq:beta-trace} and \eqref{eq:V-chain}:
\begin{equation}\label{eq:two-chains}
\Tr(M(A\otimes A))
\;=\;
\sum_{(P_1,\dots,P_N)}
\Bigl[\,\bigl(u_{\mathrm{end}}^{P_N}\bigr)^{\!\top}
W^{P_{N-1}}\cdots W^{P_3}\,u_{\mathrm{init}}^{(P_1,P_2)}\Bigr]
\Bigl[\,\beta_{P_1}^{(1)}\,\beta_{P_2}^{(2)}\cdots\beta_{P_N}^{(N)}
\bigr]\!\bigl[1,1\bigr].
\end{equation}
The two bracketed factors are independent products: the first runs over
$P_3,\dots,P_{N-1}$ in the $4$-dimensional ``replica-pattern'' space, the
second runs over $P_1,\dots,P_N$ in the $r^2$-dimensional ``MPO-bond''
space. Because the patterns at different sites are summed independently,
we can interleave the sums: at each interior site $l\in\{3,\dots,N-1\}$ the
contribution is
\[
\sum_{P_l\in\bP}\, W^{P_l}\;\otimes\;\beta_{P_l}^{(l)\,\top}.
\]
The transpose appears because the $W^{P_l}$ is acting on the replica space
in the ``$\text{out}\leftarrow\text{in}$'' direction (left multiplication),
whereas $\beta_{P_l}^{(l)}$, by definition~\eqref{eq:beta-def}, has its rows
indexed by the \emph{left} bond and its columns by the \emph{right} bond, and
thus acts on the bond space in the ``$\text{in}\rightarrow\text{out}$''
direction (right multiplication). Transposing $\beta_{P_l}^{(l)}$ aligns the
two directions so they can be combined into a single
``$\text{out}\leftarrow\text{in}$'' transfer matrix
\begin{equation}\label{eq:Wtilde-derived}
\Wtilde^{(l)} \;=\; \sum_{P\in\bP}\, W^P\,\otimes\,\bigl(\beta_P^{(l)}\bigr)^{\!\top}
\quad\in\;\IR^{4 r_l^2 \,\times\, 4 r_{l-1}^2},
\end{equation}
which matches \eqref{eq:Wtilde}.

The boundary terms are obtained by absorbing sites $l\in\{1,2\}$ into a
single initial vector and site $l=N$ into a single end vector. For the
initial vector, sites $1$ and $2$ contribute
\[
u_{\mathrm{init}}^{(P_1,P_2)}\;\otimes\;\vecop\!\bigl(\beta_{P_1}^{(1)}\,\beta_{P_2}^{(2)}\bigr),
\]
where $\vecop(\,\cdot\,)$ flattens the $1\times r_2^2$ row vector into an
element of $\IR^{r_2^2}$ (the leftmost MPO bond is closed, $r_0=1$, so
$\beta_{P_1}^{(1)}\,\beta_{P_2}^{(2)}$ is a row vector). For the end
vector, site $N$ contributes
\[
u_{\mathrm{end}}^{P_N}\;\otimes\;\vecop\!\bigl(\beta_{P_N}^{(N)}\bigr),
\]
where $\beta_{P_N}^{(N)}$ is a column vector of length $r_{N-1}^2$
(rightmost bond closed, $r_N=1$). Summing over $P_1,P_2,P_N$ gives the
expressions \eqref{eq:uinit} and \eqref{eq:uend}.

The total contraction is therefore the matrix-vector chain
\eqref{eq:thm-main}, which proves Theorem~\ref{thm:main}.

\subsection{The variance}

A small corollary, useful in applications: since $\IE[S^\top S] = I_D$,
\[
\IE\bigl[\langle S^\top S, A\rangle\bigr] \;=\; \tr(A),
\]
and consequently
\begin{equation}\label{eq:variance}
\mathrm{Var}\bigl(\langle S^\top S, A\rangle\bigr)
\;=\;
\Tr\bigl(M\,(A\otimes A)\bigr)\;-\;\tr(A)^2,
\end{equation}
where $\Tr(M(A\otimes A))$ is given by \eqref{eq:thm-main} and $\tr(A)$
is the standard MPO trace
\[
\tr(A) \;=\; \prod_{l=1}^N\Bigl(\sum_{a\in[d]} A^{(l)}_{aa}\Bigr).
\]

\section{Closed-form values of $\beta_P^{(l)}$}\label{sec:beta-explicit}

It is convenient to write each $\beta_P^{(l)}$ in terms of simpler bond
matrices. Let
\[
S_d^{(l)} \;=\; \sum_{a=1}^d A^{(l)}_{aa},
\qquad
\Delta^{(l)} \;=\; \sum_{a=1}^d A^{(l)}_{aa}\otimes A^{(l)}_{aa}.
\]
Then \eqref{eq:beta-def} unfolds as follows.
\begin{align*}
\beta_{4K}^{(l)} \;&=\; \Delta^{(l)},\\
\beta_{A}^{(l)}  \;&=\; \bigl(S_d^{(l)}\bigr)\otimes\bigl(S_d^{(l)}\bigr) \;-\; \Delta^{(l)},\\
\beta_{U_1}^{(l)} \;&=\; \sum_{a,b=1}^{d} A^{(l)}_{ab}\otimes A^{(l)}_{ab} \;-\; \Delta^{(l)},\\
\beta_{U_2}^{(l)} \;&=\; \sum_{a,b=1}^{d} A^{(l)}_{ab}\otimes A^{(l)}_{ba} \;-\; \Delta^{(l)}.
\end{align*}

\paragraph{Sanity checks.}
\begin{itemize}
\item For $A^{(l)} = I_d$ (the identity, an MPO of rank~$1$), one has
$A^{(l)}_{ab} = \delta_{ab}$, so $\Delta^{(l)} = d$ and
$\beta_{4K}^{(l)} = d$, $\beta_A^{(l)} = d^2 - d = d(d-1)$,
$\beta_{U_1}^{(l)} = \beta_{U_2}^{(l)} = 0$. These are exactly the
multiplicities used in the trace lemma of the thesis (the case $A=I$
of \eqref{eq:thm-main} reduces to $\Tr(M)$).
\item For $A^{(l)} = J_d = \mathbf{1}\mathbf{1}^\top$, one has
$A^{(l)}_{ab} = 1$ for all $a,b$, so
$\beta_{4K}^{(l)} = d$ and
$\beta_A^{(l)} = \beta_{U_1}^{(l)} = \beta_{U_2}^{(l)} = d(d-1)$.
This recovers the all-ones $A$ case discussed in the thesis.
\end{itemize}

\section{Numerical verification}\label{sec:verify}

We compute~\eqref{eq:thm-main} and compare against the brute-force
evaluation $\vecop(A)^\top M\,\vecop(A)$ for the materialized
$D^2\times D^2$ matrix $M$. Random MPO cores $A^{(l)}$ are drawn
i.i.d.~Gaussian with shape $r_{l-1}\times d\times d\times r_l$, varying
the bond rank $r$. Six configurations:

\medskip
\begin{table}[h]
\centering
\begin{tabular}{ccccccccc}
\toprule
$m$ & $N$ & $d$ & $r$ & dist.\ &
$\Tr(M(A\otimes A))$ (closed) & $\vecop(A)^\top M\,\vecop(A)$ & rel.~err.\\
\midrule
4 & 3 & 2 & 1 & gauss   & $1.20357\!\times\!10^{1}$ & $1.20357\!\times\!10^{1}$ & $8.9\!\times\!10^{-16}$\\
4 & 3 & 2 & 2 & gauss   & $2.62746\!\times\!10^{2}$ & $2.62746\!\times\!10^{2}$ & $6.5\!\times\!10^{-16}$\\
4 & 4 & 2 & 2 & rad     & $2.10802\!\times\!10^{2}$ & $2.10802\!\times\!10^{2}$ & $2.7\!\times\!10^{-16}$\\
4 & 4 & 2 & 3 & gauss   & $3.93434\!\times\!10^{4}$ & $3.93434\!\times\!10^{4}$ & $1.9\!\times\!10^{-16}$\\
4 & 5 & 2 & 2 & rad     & $1.02678\!\times\!10^{5}$ & $1.02678\!\times\!10^{5}$ & $2.8\!\times\!10^{-16}$\\
6 & 4 & 3 & 2 & gauss   & $2.68329\!\times\!10^{4}$ & $2.68329\!\times\!10^{4}$ & $4.1\!\times\!10^{-16}$\\
\bottomrule
\end{tabular}
\caption{Closed form vs.\ brute force for $\Tr(M(A\otimes A))$, agreement
to the level of double-precision rounding.}
\label{tab:verify}
\end{table}

A reference implementation in approximately 100 lines of Python is provided
in the file {\tt trace\_MB\_mpo.py}; running it as a script reproduces the
table above.

\section{Cost and remarks}

\paragraph{Computational cost.} With $r := \max_l r_l$, the dominant cost
is the propagation step \eqref{eq:Wtilde}: each $\Wtilde^{(l)}$ has size
$4 r^2 \times 4 r^2$, and its application to a vector of length $4 r^2$
costs $O(r^4)$ floating-point operations (after exploiting the Kronecker
structure, the cost is $O(r^4)$ and not $O(r^8)$). Over $N-3$ interior
sites, total cost is $O(N r^4)$. Building each $\beta_P^{(l)}$ from the
core $A^{(l)}$ costs $O(d^2 r^4)$.

\paragraph{Comparison with the row-sum bound.} The thesis row-sum lemma
gives $\Tr(M B) \le R\cdot\Tr(B)$ where $R$ is the maximum row sum of $M$.
For PSD $A$ this implies
$\mathrm{Var}(\langle S^\top S,A\rangle)\le (R-1)\,\tr(A)^2$. Numerically
$R - 1$ scales as $\alpha(d)^N/m$ with $\alpha(d)\approx d$, so the bound
is exponential in $N$. The closed form \eqref{eq:thm-main} replaces the
bound by an exact value and exhibits no such blow-up; in particular for
$A = I$ one verifies $\Tr(M)/d^{2N} \approx (1+a(d)/m)^N$ with $a(d)$ a
small constant, which is the optimal-shape regime.

\paragraph{Higher moments.} The same construction extends to
$\Tr\bigl((M(A\otimes A))^k\bigr)$ for any $k\ge 1$ by treating the $k$-fold
contraction as a chain of $k$ stacked aggregated transfer matrices, with
appropriate trace-closure boundary. The result is a matrix-vector chain on
$\IR^{(4r^2)^k}$; we omit the details here.

\appendix

\section{Reminder: explicit forms of $W^P$, $u_{\mathrm{init}}^{(P_1,P_2)}$,
$u_{\mathrm{end}}^{P_N}$}\label{app:W-formulas}

For completeness we record the formulas from the thesis. With
$\sigma_2 = 1/m$, and $\sigma_4 = \sigma_2^2$ (Rademacher) or
$\sigma_4 = 3\sigma_2^2$ (Gaussian), and $(m_1,m_2)$ the closest integer
factor pair of $m$, define for $P, Q\in\bP$ the partition meet
$P\wedge Q$ (the partition whose blocks are intersections of $P$- and
$Q$-blocks; if $P\wedge Q$ has any odd-sized block, set
$\pi(P\wedge Q):=\varnothing$ and let the moments below vanish). Define the
multiplicity
\[
N_P(M) \;=\;
\begin{cases}
M & \text{if }P = 4K,\\
M(M-1) & \text{otherwise},
\end{cases}
\]
and the moment functions
\[
\mu^X_{P,B} \;=\;
\begin{cases}
\sigma_4 & \pi(P\wedge B) = 4K,\\
\sigma_2^2 & \pi(P\wedge B) \in \{A, U_1, U_2\},\\
0 & \text{else,}
\end{cases}
\qquad
\mu^g_{L,R} \;=\;\sum_{s\in\bP} N_s(m)\,\mu^X_{L,s}\,\mu^X_{R,s}.
\]
The $4\times 4$ pattern transfer matrix is
\begin{equation}\label{eq:W-explicit}
W^P\bigl[Z_o,Z_i\bigr]
\;=\;
N_{Z_o}(m)\sum_{s_1,s_2\in\bP}
N_{s_1}(m_1)\,N_{s_2}(m_2)\,
\mu^X_{P,\,s_1\wedge s_2}\,
\mu^g_{Z_i\wedge s_1,\,s_2\wedge Z_o},
\end{equation}
and the boundary vectors $u_{\mathrm{init}}^{(P_1,P_2)},u_{\mathrm{end}}^{P_N}$
are given by analogous explicit polynomials (formulas (20), (21) of the
thesis). These quantities depend only on $m$ and the noise distribution,
not on $A$, $N$, or $d$.

\end{document}
